\chapter{Asset Management}

The Asset Management subsystem in the Caffeine Engine is designed around the principles of asynchronous resource loading, zero-copy memory access, and automated lifetime management through reference counting. By decoupling the acquisition of a resource from its physical residence in memory, the engine maintains high frame rates during heavy I/O operations and ensures that memory fragmentation is minimized through the use of dedicated linear allocators for each asset.

\needspace{6\baselineskip}
\section{The Asset Handle}

The primary interface for interacting with the Asset Manager is the \texttt{AssetHandle<T>}. This template class serves as a RAII-compliant reference-counted handle to a resource of type $T$.

\subsection{Reference Counting Semantics}

\texttt{AssetHandle<T>} manages the lifetime of an asset by communicating with the \texttt{AssetManager}. The reference count is incremented whenever a handle is copied and decremented when it is destroyed.

\begin{equation}
    R_{current} = \sum_{i=1}^{n} H_i
\end{equation}

Where $R_{current}$ is the total reference count of an asset and $n$ is the number of active \texttt{AssetHandle} instances pointing to that specific resource ID. When $R_{current}$ drops to zero, the asset becomes a candidate for eviction during the next garbage collection cycle.

\needspace{5\baselineskip}
\subsection{Implementation Details}

The handle is designed to be lightweight, containing only a pointer to the manager and the unique 32-bit identifier for the asset.

\begin{lstlisting}[language=C++, caption={AssetHandle interface}]
template<typename T>
class AssetHandle {
public:
    AssetHandle();
    AssetHandle(AssetManager* mgr, u32 id);
    ~AssetHandle();

    AssetHandle(const AssetHandle& o);
    AssetHandle& operator=(const AssetHandle& o);

    AssetHandle(AssetHandle&& o) noexcept;
    AssetHandle& operator=(AssetHandle&& o) noexcept;

    bool     isValid() const;
    bool     isReady() const;
    const T* get()     const;

    explicit operator bool() const;
    u32 id() const;
};
\end{lstlisting}

The move constructor and move assignment operator are specifically optimized to transfer ownership without triggering atomic increments or decrements in the \texttt{AssetManager}. After a move, the source handle is set to a null state (\texttt{m_id = ~0u}), ensuring that the destructor does not decrement the count for the transferred resource.

\begin{specbox}{Design Rationale: Handle-Based Access}
By returning a handle instead of a raw pointer, the engine can safely perform hot-reloading or garbage collection in the background. The user never holds a direct pointer to the underlying data for longer than a single frame's scope, allowing the manager to invalidate or move memory without causing dangling pointers in game logic.
\end{specbox}

\needspace{6\baselineskip}
\section{The Asset Manager}

The \texttt{AssetManager} is the central authority for resource lifecycle. It maintains an internal registry of assets, indexed by their path, and coordinates with the \texttt{JobSystem} for background loading.

\subsection{Loading Paths}

The manager provides two distinct paths for loading resources: synchronous and asynchronous.

\begin{itemize}
    \item \textbf{Sync Path:} Invoked via \texttt{loadSync<T>()}. This method calls \texttt{loadInternal()} immediately, blocking the calling thread until the asset is fully loaded and resolved. This is typically used during initial engine startup or loading screens where immediate availability is required.
    \item \textbf{Async Path:} Invoked via \texttt{loadAsync<T>()}. This method registers the asset requirement and calls \texttt{scheduleLoad()}, which submits a task to the \texttt{JobSystem}. The calling thread receives a handle immediately, but \texttt{isReady()} will return false until the background task completes.
\end{itemize}

\needspace{5\baselineskip}
\subsection{Zero-Copy Memory Model}

One of the most performance-critical features of the Caffeine Engine is its zero-copy asset storage. Each \texttt{AssetEntry} contains a \texttt{std::unique_ptr<LinearAllocator>} that owns the memory slab where the raw file data is loaded.

\begin{equation}
    M_{total} = \sum_{j=1}^{m} S_j + \Omega
\end{equation}

Where $M_{total}$ is the total memory footprint, $S_j$ is the size of the allocator slab for asset $j$, and $\Omega$ represents the fixed overhead of the manager's tracking structures.

When an asset is loaded, the engine maps its structures directly onto the memory-mapped buffer. For example, a \texttt{Texture} object does not contain a copy of the pixel data; instead, it contains a pointer that points directly into the \texttt{LinearAllocator} slab.

\begin{lstlisting}[language=C++, caption={AssetEntry Structure}]
struct AssetEntry {
    std::string                      path;
    AssetType                        cafType;
    std::atomic<LoadStatus>          status;
    std::unique_ptr<LinearAllocator> allocator;
    const void*                      payload;
    ResolvedData                     resolved;
    std::atomic<u32>                 refCount;
    u64                              sizeBytes;
};
\end{lstlisting}

\needspace{5\baselineskip}
\subsection{Garbage Collection and Cache Management}

The \texttt{collectGarbage()} function is responsible for pruning the cache. It iterates through the asset registry and identifies entries where the \texttt{refCount} is zero. These entries are evicted, and their associated \texttt{LinearAllocator} slabs are freed, returning memory to the system.

The manager also tracks performance metrics through the \texttt{CacheStats} structure:

\begin{lstlisting}[language=C++, caption={CacheStats definition}]
struct CacheStats {
    u64 totalCachedBytes;
    u64 maxCacheBytes;
    u32 textureCount;
    u32 audioCount;
    u32 pendingJobs;
    f32 cacheHitRate;
};
\end{lstlisting}

The hit rate is calculated as:
\begin{equation}
    HitRate = \frac{C_{hits}}{C_{total}}
\end{equation}
where $C_{hits}$ is the number of times an already-loaded asset was requested via path lookup, and $C_{total}$ is the total number of load requests.

\needspace{6\baselineskip}
\section{Asset Runtime Types}

Caffeine utilizes a specialized \texttt{.caf} format for its assets. At runtime, these are represented as thin views into the loaded data.

\subsection{Asset Mapping with Type Traits}

To support a generic template-based API, the engine uses the \texttt{AssetTypeTrait<T>} mechanism. This allows the \texttt{AssetManager} to determine the internal \texttt{AssetType} discriminator at compile time.

\begin{lstlisting}[language=C++, caption={AssetTypeTrait Specialization}]
template<> struct AssetTypeTrait<Texture> {
    static constexpr AssetType cafType = AssetType::Texture;
};
\end{lstlisting}

\needspace{5\baselineskip}
\subsection{Specific View Structures}

The three primary asset types currently supported by the engine are \texttt{Texture}, \texttt{AudioClip}, and \texttt{ShaderBlob}.

\begin{itemize}
    \item \textbf{Texture:} Contains dimensions, format, and a pointer to the pixel buffer.
    \begin{lstlisting}[language=C++]
struct Texture {
    u32       width;
    u32       height;
    u32       format;
    u32       mipLevels;
    const u8* pixels;
    u64       pixelDataSize;
};
    \end{lstlisting}
    
    \item \textbf{AudioClip:} Holds PCM data views and metadata for audio playback.
    \begin{lstlisting}[language=C++]
struct AudioClip {
    u32       sampleRate;
    u16       channels;
    u16       bitsPerSample;
    u32       sampleCount;
    const u8* pcmData;
    u64       pcmDataSize;
};
    \end{lstlisting}
    
    \item \textbf{ShaderBlob:} A view into compiled shader bytecode.
    \begin{lstlisting}[language=C++]
struct ShaderBlob {
    u32       stage;
    const u8* bytecode;
    u64       bytecodeSize;
};
    \end{lstlisting}
\end{itemize}

\begin{specbox}{Optimization: Alignment and Padding}
When resolving these views, the \texttt{AssetManager} ensures that pointers (\texttt{pixels}, \texttt{pcmData}, \texttt{bytecode}) are aligned to the requirements of the underlying hardware (e.g., SIMD alignment for audio or GPU alignment for textures), even though they reside within a shared linear buffer.
\end{specbox}
